\documentclass{article}
\usepackage{amsmath, amssymb}

\title{lec511}
\author{Dylan Tang}
\date{May 2026}

\begin{document}

\maketitle

\section{Principal Component Analysis}
\begin{itemize}
    \item PCA finds a new, lower-dimensional basis that attempts to capture as much variance in the data as possible
        \item Recall we can find PCA of rank k with SVD.
        \[ 
            X_k = U_k\Sigma V_k^T
        \]
    \item PCA is useful in that the degrees of freedom of plausible values is often much larger than the degrees of freedom that actually make sense in the real world.
    \item \textbf{Information Bottleneck Theory} - All significant learning requires setting an information bottleneck
        \begin{itemize}
            \item Information bottleneck - real-world data often lies in some lower dimensional space when compared to the entire subset of possible data values
        \end{itemize}
\end{itemize}
\section{Neural Networks}
\subsubsection{Forward Propagation}
\begin{itemize}
    \item Consider input observations $X_1, \dots, X_n$, where each $X_i \in \mathbb{R}^p$.
    
    \item At each hidden layer, compute a weighted linear combination of the inputs, then apply a nonlinear activation function:
    \[
        a_i^{(j)} = \sigma\left(\left[w^{(j)}\right]^T X_i + b^{(j)}\right),
    \]
    where $w^{(j)}$ is the weight vector for hidden unit $j$, $b^{(j)}$ is the bias term, and $\sigma(\cdot)$ is the activation function.
    
    \item More generally, for layer $\ell$, the forward propagation step is
    \[
        Z^{(\ell)} = W^{(\ell)}A^{(\ell - 1)} + b^{(\ell)},
    \]
    \[
        A^{(\ell)} = \sigma\left(Z^{(\ell - 1)}\right),
    \]
    where $A^{(0)} = X$ is the input layer.
    
    \item The final layer produces the network output:
    \[
        \hat{Y} = A^{(L)}.
    \]
    \item At the final layer, the network produces a prediction $\hat{Y}_i$ for each input $X_i$:
    \[
        \hat{Y}_i = f(X_i; \theta),
    \]
where $\theta$ represents all weights and biases in the network.
\end{itemize}

\subsubsection{Gradient Descent}
Gradient descent is given by:
\[
    W^{(t+1)}
    =
    W^{(t)}
    -
    \eta \nabla_{W}\mathcal{L}\left(W^{(t)}\right)
\]
\subsubsection{Backward Propogation}
\begin{itemize}
\item For a hidden layer $\ell$, recall that
\[
    Z^{(\ell)} = W^{(\ell)}A^{(\ell - 1)} + b^{(\ell)},
    \qquad
    A^{(\ell)} = \sigma\left(Z^{(\ell)}\right).
\]
\end{itemize}
Using the chain rule, the gradient with respect to $W^{(\ell)}$ is
\[
    \frac{\partial \mathcal{L}}{\partial W^{(\ell)}}
    =
    \frac{\partial \mathcal{L}}{\partial A^{(\ell)}}
    \frac{\partial A^{(\ell)}}{\partial Z^{(\ell)}}
    \frac{\partial Z^{(\ell)}}{\partial W^{(\ell)}}.
\]

Since
\[
    \frac{\partial A^{(\ell)}}{\partial Z^{(\ell)}}
    =
    \sigma'\left(Z^{(\ell)}\right),
    \qquad
    \frac{\partial Z^{(\ell)}}{\partial W^{(\ell)}}
    =
    \left(A^{(\ell - 1)}\right)^T,
\]
we get
\[
    \boxed{
    \frac{\partial \mathcal{L}}{\partial W^{(\ell)}}
    =
    \left[
    \frac{\partial \mathcal{L}}{\partial A^{(\ell)}}
    \odot
    \sigma'\left(Z^{(\ell)}\right)
    \right]
    \left(A^{(\ell - 1)}\right)^T
    }
\]
Recall that the recurrence relation is  \[
    A^{(\ell)}
    =
    \sigma\left(W^{(\ell)}A^{(\ell-1)} + b^{(\ell)}\right)
\] so we could find $ \frac{\partial \mathcal{L}}{\partial A^{(\ell)}}$ by repeatedly  unwrapping the recurrence relation.

\textbf{Remarks}
\begin{itemize}
    \item Gradient descent requires optimizing for each weight connecting nodes from $\ell$ to $\ell + 1$ - expensive
    \item It is important to standardize $X_1 \dots X_N$ to prevent exploding/vanishing gradients
\end{itemize}
\subsubsection{Code Snippets}
\textbf{Standard Regression using Neural Nets}
\begin{verbatim}
    model = tf.keras.models.Sequential() // standard layer-by-layer NN
    model.add(tf.keras.layers.InputLayer(input_shape=(8,))) // X1...X8
    model.add(tf.keras.layers.Dense(units=1)) // one hidden layer with no activation

    model.compile(
        optimizer=tf.keras.optimizers.Adam(),
        loss=tf.keras.losses.MeanSquaredError()
    )

    history = model.fit(
        x=x,
        y=y,
        batch_size=1000,
        epochs=epochs,
        shuffle=True
    )
\end{verbatim}

\textbf{2 Dense Layers, ReLu activation}
\begin{verbatim}
    model = tf.keras.models.Sequential()

    model.add(tf.keras.layers.InputLayer(input_shape=(8,)))

    model.add(tf.keras.layers.Dense(units=16, activation='relu'))
    model.add(tf.keras.layers.Dense(units=8, activation='relu'))

    model.add(tf.keras.layers.Dense(units=1))

    model.compile(
        optimizer=tf.keras.optimizers.Adam(),
        loss=tf.keras.losses.MeanSquaredError()
    )

    history = model.fit(
        x=x,
        y=y,
        batch_size=1000,
        epochs=epochs,
        shuffle=True
    )
\end{verbatim}
\begin{itemize}
    \item \textbf{Adam} optimizatin is gradient descent with momentum and adaptive learning rates.
\end{itemize}
\textbf{Convolutional Neural Network}
\begin{itemize}
    \item Instead of having only activation functions as bottlenecks, has more complex hidden layers with pooling, stride, padding, and filter.
\end{itemize}
\begin{verbatim}
import tensorflow as tf

model = tf.keras.models.Sequential()

# Input image: 28 x 28 grayscale image
model.add(tf.keras.layers.InputLayer(input_shape=(28, 28, 1)))

# First convolution block
model.add(tf.keras.layers.Conv2D(
    filters=32,
    kernel_size=(3, 3),
    strides=(1, 1),
    padding='same',
    activation='relu'
))

# Pooling layer
model.add(tf.keras.layers.MaxPooling2D(
    pool_size=(2, 2),
    strides=(2, 2)
))

# Second convolution block
model.add(tf.keras.layers.Conv2D(
    filters=64,
    kernel_size=(3, 3),
    strides=(1, 1),
    padding='same',
    activation='relu'
))

# Pooling layer
model.add(tf.keras.layers.MaxPooling2D(
    pool_size=(2, 2),
    strides=(2, 2)
))

# Flatten feature maps into vector
model.add(tf.keras.layers.Flatten())

# Fully connected layer
model.add(tf.keras.layers.Dense(
    units=128,
    activation='relu'
))

# Output layer for 10-class classification with softmax loss
model.add(tf.keras.layers.Dense(
    units=10,
    activation='softmax'
))

model.compile(
    optimizer=tf.keras.optimizers.Adam(),
    loss=tf.keras.losses.SparseCategoricalCrossentropy(),
    metrics=['accuracy']
)

history = model.fit(
    x=x_train,
    y=y_train,
    batch_size=64,
    epochs=epochs,
    shuffle=True,
    validation_data=(x_test, y_test)
)
\end{verbatim}
\end{document}